\cotefile{157-1}
\scanpage{34}\mspage{32}
\uline{Dér} 4 : Le morphisme de chgt. de base pour images directes \inserted{$f_*$ \deleted{\unsure{rel.}}} associé à un diagramme (1) est toujours un iso (dérivateurs à g). Condition duale pour dérivateurs à droite et le chgt de base pour images directes $f_!$, pour le diagramme
\diagram{$(1')$\qquad\begin{tikzcd}[ampersand replacement=\&,column sep=large]
 {}_{j_0\backslash}I \arrow[r] \arrow[d] \& I \arrow[d,"f",""{name=right,coordinate}] \\
 e \arrow[r,"p=\alpha_{j_0}"',""{name=bottom,coordinate}] \& J
 \arrow[from=bottom,to=right,shorten <=2pt,shorten >=2pt,bend left=10,"\unsure{\gamma|r}"]
 \end{tikzcd}}
associé à $f:I\to J$ dans Diag, et $j_0\in\mathrm{Ob}\,J$.

% slanted
\margin{\uline{Totalement} \unsure{canulé}, mais \unsure{raisonnable|raison} \unsure{aussi}, pour coeff \unsure{constants} sur $\unsure{I|F}$ \unsure{Version} \unsure{hot}}
\unsure{Dér} 5 : Soit $f:I\to J$ une \inserted{\unsure{hot-|hoto-}}\uline{équivalence} \uline{stricte}, i.e. $\forall\,j_0\in\mathrm{Ob}\,J$, $I_{/j_0}$ est \unsure{contractiles|contractibles}. Alors $f^*$ est pl. fidèle, i.e. le morphisme d'adjonction
\[
 \mathrm{id}_{\mathbb D(J)}\longrightarrow f_*f^*
\]
est un iso ($\mathbb D$ dérivateur à g.). Dualemt, si $f$ est \deleted{\unsure{tel que} $f^\circ : I^\circ$} \inserted{une équivalence} co-stricte, i.e. $f^\circ:I^\circ\to J^\circ$ est une \unsure{équiv.|équ.} stricte, i.e. $\forall\,j_0\in\mathrm{Ob}\,J$, ${}_{j_0\backslash}I$ est \unsure{contractile|contractible}, alors (si $\mathbb D$ est un dérivateur à dr.) le foncteur $f^*$ est pl. \unsure{fidèle|fidèl}, i.e. le foncteur
\[
 f_!f^*\longrightarrow\mathrm{id}_{\mathbb D(J)}
\]
\unsure{est} \deleted{\illegible{1}} un iso.

\uline{NB} Donc les foncteurs $f_*$ (resp $f_!$) sont des foncteurs de localisation.
% archive sheet number: 33
\scanpage{35}\mspage{41}
% archive sheet number: 34
\deleted{Donc si
\[
 \operatorname{hom}(e,F)\longrightarrow\operatorname{hom}(X,F)
\]
est iso pour tout $F$ dans $\mathrm{Hot}$, \unsure{passant|prenant} aux $\pi_0$ on en déduit que
\[
 \operatorname{Hom}_{\mathrm{Hot}}(e,F)\longrightarrow\operatorname{Hom}_{\mathrm{Hot}}(X,F)
\]
est iso pour \unsure{tout} $F$, donc $X\to e$ iso dans $\mathrm{Hot}$, i.e. $X$ contractile, OK.

Mais le cas général \unsure{de la|du} ``\unsure{propr.|principe}'' reste douteux \ldots}

\margin{\uline{NB} les pages 33 à 40 ont été supprimées, car \unsure{canulées}}

Je vais reprendre la terminologie et les axiomes relatifs aux dérivateurs.

Un \uline{prédérivateur} est un 2-foncteur
\[
 \mathrm{Diag}^{\circ}\xrightarrow{\mathbb D}(\mathrm{Cat}),
\]
% slanted
\margin{Diag : \uline{domaine} du dérivateur. Les objets sont les \uline{cat. d'}\deleted{\unsure{coefficients}}\inserted{indices} pour $\mathbb D$. Les $\mathbb D(I)$ sont les cat. de \unsure{coeffs} pour $\mathbb D$, et les objets de $\mathbb D(I)$ \unsure{des} ``\unsure{coefficients|coeffts} sur $I$''}
$(\mathrm{Diag})$ est une \inserted{2-}sous-catégorie pleine de $(\mathrm{Cat})$. On \unsure{la supposera|suppose} tout au moins que Diag contient les \uline{ens.} \uline{ordonnés} \uline{finis}, que Diag est stable par $\underleftarrow{\lim}$ finies, par sommes finies. Qu'avec toute catégorie $C$ et toute sous-catégorie strictemt pleine ouverte (resp. fermée) elle contienne la sous-catégorie complémentaire.
\scanpage{36}\mspage{42}
% slanted
\margin{Notation $f^*$ : Image inverse de coefficients.}
On définit le prédérivateur \inserted{$\mathbb D_{\unsure{\mathcal M|M},\Sigma}$} associé à un système
\[
 (\mathcal M,\Sigma)\qquad \mathcal M\unsure{\in}\ \text{Catégorie},\quad\Sigma\subset\unsure{\mathrm{FC}|\mathrm{Fl}}(\mathcal M)
\]
Si $\Sigma=\varnothing$, on \deleted{\unsure{pose}} écrit $\mathbb D_{\mathcal M}$

Les prédérivateurs forment une \inserted{2-}catégorie\inserted{ $\mathrm{PREDER}$}, et
\[
 (\mathcal M,\Sigma)\longmapsto\deleted{\unsure{\mathbb D}}\mathbb D_{\mathcal M,\Sigma}
\]
est un 2-foncteur
\[
 \mathrm{PREMOD}\left(\begin{gathered}\text{``Cat \inserted{des Catégories} de modèles''}\\ \text{(\unsure{faible|faibles})}\end{gathered}\right)\longrightarrow\mathrm{PREDER}
\]

\inserted{Pour \unsure{$D$|$\mathbb D$} donné,} On définit les dérivateurs induits ($I\in\unsure{\mathrm{Diag}}$)
\[
 \mathbb D_I=(J\longmapsto\mathbb D(I\times J)),
\]
\unsure{dépendant|dépendent} de $I$ de façon contravariante, d'où un 2-\deleted{\unsure{contra}}foncteur
\[
 \begin{aligned}
 \mathrm{Diag}^{\circ}&\longrightarrow\mathrm{PREDER}\\
 I&\to\mathbb D_I
 \end{aligned}
\]
ou même un bi-2 foncteur
\[
 \begin{aligned}
 \mathrm{PREDER}\times\mathrm{Diag}^{\circ}&\longrightarrow\mathrm{PREDER}\\
 (\mathbb D,I)&\to\mathbb D_I
 \end{aligned}
\]

Si $\mathbb D$ est défini par $(\mathcal M,\Sigma)$, alors $\mathbb D_I$ est défini par $\mathcal M(I),\Sigma(I)$ :
\[
 (\mathbb D_{\mathcal M,\Sigma})_I\simeq\mathbb D_{\mathcal M(I),\Sigma(I)},
\]
% archive sheet number: 35
\scanpage{37}\mspage{43}
% archive sheet number: 36
où $\unsure{\mathcal M|M}(I)=\underline{\mathrm{Hom}}(I^\circ,\mathcal M)$ (\unsure{préfs|préf.} sur $I$ à val. dans $\mathcal M$), et $\Sigma_{\unsure{I|J}}\subset\mathrm{Fl}\,\mathcal M(I)$ est déduit de $\Sigma$ en regardant fibre par fibre (supposant maintenant $\Sigma$ saturé).

Pour tout axiome qu'on \unsure{dégagera|dégage} pour certains prédérivateurs, il faudra vérifier qu'ils sont stables par passage \unsure{au|aux} \deleted{\unsure{dé}}prédérivateurs induits $\mathbb D_I$, et aussi \deleted{\unsure{pour que}} comment ils se comportent par dualité (passage de $\mathbb D$ à $\mathbb D^\circ$).

Définition du prédér. dual $\mathbb D^\circ$
\[
 \mathbb D^\circ(I)\overset{\mathrm{def}}{=}\mathbb D(I^\circ)^\circ\qquad\text{etc.}
\]

% slanted
\margin{Premier Axiome de localisation}
\uline{Der 1} \quad a) Si $I,J\in\mathrm{Diag}$, alors
\[
 \mathbb D(I\amalg J)\longrightarrow\mathbb D(I)\times\mathbb D(J)
\]
est une équivalence de catégories

b) $\mathbb D(\varnothing)$ est la catégorie finale

Stable par passage aux \deleted{\unsure{dé}}prédér. induits, et passage au dual.

\uline{Cor} \quad Soit $(I_\alpha)_{\alpha\unsure{\in|+}A}$ famille finie de cat $\in\mathrm{Diag}$ $I=\coprod_{\alpha\in A}I_\alpha$. Alors
\[
 \mathbb D(I)\longrightarrow\prod_{\alpha\in A}\mathbb D(I_{\unsure{-|\alpha}})
\]
une équivalence de catégories.

% slanted
\margin{deuxième axiome de localisation. Stable par induction et par dualité}
\uline{Der 2} \quad \inserted{$\forall I\in\mathrm{Diag}$} le foncteur canonique
\[
 \deleted{\unsure{I}}\mathbb D(I)\longrightarrow\underline{\mathrm{Hom}}(I^\circ,A_{\mathbb D})
\]
\margin{$=\mathbb D(e)$ catégorie fondamentale pour le dérivateur}
est conservatif
\scanpage{38}\mspage{44}
Un prédérivateur est appelé un \uline{dérivateur à gauche} s'il satisfait Der 1\inserted{ et \unsure{Der|Dér} 2}, et de plus \deleted{\unsure{l'axiome suivant}} les axiomes suivants (dont voici le premier).

% slanted
\margin{Existence des images directes (stable par induction\unsure{.}}
\uline{Der 3} (g) : Pour toute $f:I\to J$ dans Diag, le foncteur $f^*:\mathbb D(J)\to\mathbb D(I)$ admet un \uline{adjoint à droite} $f_*$.

On appelle $f_*$ le foncteur image directe \inserted{(cohomologique)} pour $f$. L'\unsure{objet|obj} $f_*(F)\inserted{{}\in\mathbb D(J)}$ (\inserted{pour }$F\in\mathbb D(I)$) s'appellera aussi \deleted{de} \inserted{l'objet} ``\uline{cohomologie relative}'' de $I$ sur $J$ (par $f$), à \unsure{coëfficients|coéfficients} dans $F$. On a des iso. \unsure{can|can.}
\[
 (gf)_*\simeq g_*f_*
\]

\uline{Cor.} Les catégories $\mathbb D(I)$ ont des produits finis, les foncteurs $f^*$ et $f_*$ y commutent.

Utilise Der 1 \deleted{\unsure{et D}} pour $I=\varnothing$ et $\inserted{I={}}\unsure{.}e\amalg e$, et Der 3 pour $I\xrightarrow{f}e$.

On a un \deleted{\unsure{foncteur}} \inserted{axiome} dual

\uline{Der 3} (d) : existence des adj. à g. $f_!$ des $f^*$

le foncteur $f_!$ s'appelle ``image directe homologique'' pour $f$, et on a une notion d'\uline{homologie relative} de $I$ sur $J$ (par $f$), à coeff. dans $F$.

Par la suite, la mention (g) ou (d) à côté de la désignation Der $n$ d'un axiome, signifie qu'il s'applique aux dériv. à g., resp. \unsure{à} dr. De \unsure{=|même} pour les énoncés, définitions etc.
% archive sheet number: 37
\scanpage{39}\mspage{45}
\deleted{Un prédérivateur est dit}
% slanted
\margin{Inutile comme axiome, \unsure{va|car} \unsure{devenir|devient} un théorème (résulte des autres axiomes}
\uline{Der 4}$_0$ (g) : Soit
\diagram{\begin{tikzcd}[ampersand replacement=\&]
 U \arrow[r,hook,"j"'] \& X \& Y \arrow[l,hook',"i"]
\end{tikzcd}}
\inserted{dans Diag}, \unsure{imm|immersions} ouverte et fermée complémentaires l'une de l'autre. Alors le morphisme d'adjonction dans $\mathbb D(Y)$
\[
 i^*i_*(F_Y)\longrightarrow F_Y
\]
est \deleted{toujours \unsure{plein}} \inserted{un iso}, donc \deleted{\unsure{$i_*$ est}}
\[
 i_*:\mathbb D(Y)\lhook\joinrel\longrightarrow\mathbb D(X)
\]
est pleinement fidèle, et $i^*$\inserted{$:\mathbb D(X)\to\unsure{\mathbb D(Y)|\mathbb D Y}$} est un foncteur de localisation. \deleted{Plus précisément, pour \inserted{$u:F\to G$} dans \unsure{$I$|$J$}, \deleted{\unsure{si $u$}} \unsure{$u$} est un iso ssi $\unsure{i|i^*}$ $i^*(u):i^*(F)\to i^*(G)$ \unsure{est un iso}} L'image essentielle \inserted{de $i_*$} est formée des objets $F$ de $\mathbb D(X)$ tels que $j^*(F)$ soit l'objet final de $\mathbb D(U)$.

\uline{Der \unsure{4|5}}$_0$ (d) Avec les notations précédentes pour $i,j$, le morphisme d'adjonction dans $\mathbb D(\unsure{u|U})$
\[
 F_U\to j_{\deleted{!}}^*j_{\unsure{!}}^{\deleted{\unsure{*}}}F_{\unsure{u|U}}\ \deleted{\unsure{\to F_U}}
\]
est un iso -- $j_{\unsure{*|!}}$ est pl. fid, $j^*$ un foncteur localisation. L'image essentielle de $j_!$ est la sous-catégorie pleine de $\mathbb D(X)$ formée des $F$ \deleted{\unsure{de}} tels que $i^*(F)$ soit \deleted{\unsure{l'}}\unsure{obj} initial \deleted{\unsure{de}} de $\mathbb D(Y)$.

Il y a les deux carrés de la page 8 et 9, mais pour interpréter \unsure{les|le} \unsure{caractères|caractère|condition} cartésien ou cocartésien c'est prématuré ici.
% archive sheet number: 38
\scanpage{40}\mspage{46}
\unsure{Axiomes|Axiome} de changement de base

% slanted
\margin{Fibres d'une image directe}
\uline{Der \deleted{5}4} (g) Soient $f:X\to Y$ dans Diag, $y\in Y$, considérons le diagramme \unsure{comm.|comm}
\diagram{\begin{tikzcd}[ampersand replacement=\&,column sep=large]
 X_{/y} \arrow[r,"q"] \arrow[d,"f'"'] \& X \arrow[d,"f"]\\
 e \arrow[r,"p=\alpha_y"'] \& Y
\end{tikzcd}}
et le morphisme de chgt. de base correspondant
\[
 \begin{aligned}
 &\deleted{\unsure{f_*(F)_y^*f}}\\
 &\underbrace{p^*f_*(F)}_{\substack{\parallel\,\mathrm{def}\\ f_*(F)_y}}\longrightarrow
 f'_* (q^*(F))\left(\overset{\mathrm{def}}{=}\Gamma_{\mathbb D}(X_{/y},q^*(F))\right)
 \end{aligned}
\]
C'est toujours un iso.

% slanted
\margin{Axiome d'asphéricité (coeff. constants) Mais est contenu dans \uline{Dér 6}}
\deleted{\deleted{\unsure{Dér 6}}\ \deleted{\unsure{Cor.}}\uline{Der 6} (g) : Soit $g:Z\to e$ dans Diag, et $M\in\mathbb D(e)=A$ (cat fondamentale). Supposons $\unsure{S|Z}$ contractile. Alors le morphisme d'adjonction
\[
 M\longrightarrow g_*g^*M.\quad\left(\overset{\mathrm{def}}{=}\Gamma_{\mathbb D}(Z,g^*M)\right)
\]
\margin{$\mathbb D$-coh. ``globale'', -- ``absolue, de $Z$ à coeff. dans $g^*(M)$}
est un iso.

\uline{Cor. 1} Soit $f:X\to Y$ qui soit une \inserted{\unsure{hoto-|hot-}}équiv. stricte, i.e. $\forall y\in Y$, $X_{/y}$ contractile. Alors pour $F\in\mathbb D(Y)$, \unsure{le|la} \unsure{hom|flèche} d'adjonction
\[
 F\to f_*f^*(F)
\]
est un iso.}
% archive sheet number: 39
\scanpage{41}\mspage{47}
% slanted
\margin{\deleted{\unsure{Principe}}\inserted{Axiome} de cofinalité \deleted{cohomologique} pour la cohomologie Il est \unsure{bcp plus|trop} \unsure{spécial|faible} \unsure{d'ailleurs} \unsure{faible !}}
\uline{Der 5} (g) : Soit $\alpha\deleted{\unsure{f}}:X'\to X$ dans Diag, ``cohomologiquement cofinal'' les ${}_{x\backslash}X'$ sont contractiles ($x\in X$) -- i.e. \inserted{$\alpha^\circ:X'^\circ\to X^\circ$} \unsure{est une} hoto-équivalence ``loc. sur $X^\circ$''. Alors $\forall F\in\mathbb D(X)$, l'hom
\[
 \Gamma_{\mathbb D}(X,F)\longrightarrow\Gamma_{\mathbb D}(X',\alpha^*F)
\]
\margin{\uline{NB} Si $f:X\to e$, on pose $\Gamma_{\mathbb D}(X,F)=f_*(F)$ (coh. ``absolue'' de $X$, à coeff dans $F$)}
est un iso (dans $A=\mathbb D(e)$). Même \deleted{chose si $\alpha$ est ``homologiquement cofinal'' i.e. les $X'_{/x}$ sont contractiles.}
\margin{Faux ! $\longrightarrow$}

\uline{Remarque} On devrait pouvoir renforcer considérablement cet axiome, en exigeant que
\[
 \alpha_!(\varepsilon_{\unsure{X|x}'})\xrightarrow{\sim}\varepsilon_X\qquad\text{dans }\mathrm{HOT}(X),
\]
où $\varepsilon_{X'},\varepsilon_X$ les objets finaux de $\mathrm{HOT}(X),\mathrm{HOT}(X')$, et où $\alpha_!(\varepsilon_{X'})\to\varepsilon_X$ est l'hom d'adjonction, compte tenu que $\varepsilon_{X'}\simeq\alpha^*(\varepsilon_X)$. On dit que $X'$ est \uline{homotopiquement trivial} \uline{sur} $X$. (Et si la condition duale est satisfaite dans $\mathrm{HOT}(X^\circ)$, on dit que $X'$ est cohomotopiquement trivial sur $X$.) Cette condition n'est sans doute \uline{pas} autoduale. Elle est vérifiée aussi bien si $\alpha$ est coh. cofinal, que s'il est \unsure{homologiquement|homotopiquement} cofinal (\inserted{\unsure{hot-|hoto-}}équiv \unsure{costricte|co-stricte} ou stricte)

\inserted{\uline{NB} Pour $X\xrightarrow{\unsure{f}}e$, dire que $f$ est coh. cofinal, hom cofinal, ou que $X$ est contractile, sont équivalents.}

% slanted
\margin{Th. d'asphéricité pour les coeff. constants, en cohomologie}
\uline{Cor. 1} (g) Soit $f:\deleted{\unsure{Z}}X\to e$ dans Diag, \inserted{$X$ contractile,} $M\in\mathbb D(e)$. Alors l'hom d'adjonction
\[
 M\xrightarrow{\sim}f_*f^*(M)\quad\left(\overset{\mathrm{def}}{=}\Gamma(\deleted{\unsure{Z}}X,f^*M)\right)
\]
est un iso.

% slanted
\margin{\uline{Faux} \unsure{sauf} pour coeff. constants}
\deleted{\uline{Cor. 2} (g) : Soit $f:X\to Y$ hot-équiv stricte (i.e. $\forall y\in Y$, $X_{/y}$ contractile). Alors pour $F\in\mathrm{Ob}\,\mathbb D(Y)$
\[
 F\xrightarrow[\sim]{\mathrm{adj.}}f_*f^*(F)\qquad\text{iso}
\]}
% archive sheet number: 40
\scanpage{42}\mspage{48}
On applique Dér 2, Dér 4 (g), pour se ramener à regarder, \inserted{pour $y_0\in Y$}, l'hom $\alpha$ du diagramme \inserted{(commutatif !)} suivant :
\diagram{\begin{tikzcd}[ampersand replacement=\&,column sep=large]
 \Gamma(\{y_0\},F_{y_0})\unsure{=|\leftarrow|\simeq}F_{y_0}\arrow[r,"\alpha"] \& \Gamma_{\mathbb D}(X_{/y_0},f^*(F))\\
 \Gamma(Y_{/y_0},F)\arrow[u,leftrightarrow,"\varphi","\unsure{\sim|?}"']\arrow[ur,"\alpha'"'] \&
\end{tikzcd}}
\margin{$\longleftarrow$ \deleted{défini par le foncteur $\unsure{\{y\}|f_y}$}}
où $\alpha'$ est déduit du foncteur induit par $f$
\[
 X_{/y_0}\longrightarrow Y_{/y_0}
\]
et où $\varphi$ est déduit du foncteur
\[
 \{y_0\}\xrightarrow[\substack{(\text{objet}\\ \text{final})}]{} Y_{/y_0}
\]

% slanted
\margin{Cf cor. 3 page 50, qui devrait remonter ici.}
Donc $\varphi$ est un iso par Der 5 (g)$^*$ (NB Si $z\in Z$ d'où $\{z\}\to Z$, ce foncteur est coh. cofinal ssi $z$ est objet final, \deleted{\unsure{donc $f$}} \unsure{hom.|homo} final ssi $z$ est objet initial) Donc la conclusion du corollaire équivaut à dire que $\alpha'$ est iso. Mais $X_{/y_0}\to Y_{/y_0}$ est aussi hom. cofinal \deleted{\unsure{\(\longrightarrow\)}} ; car si $y_1\xrightarrow{u}y_0$ \deleted{\unsure{objet}} $\in$ de $Y_{/y_0}$, on a un iso. can.
\[
 (X_{/y_0})_{/(y_1,u)}\simeq X_{/y_1}\qquad\text{qui est contractile}
\]
Donc par la \uline{deuxième partie} de \deleted{Dér} (5)(g), on gagne.
% slanted
\margin{qui est \uline{fausse}}
% archive sheet number: 41
\scanpage{43}\mspage{49}
On applique \uline{Der \deleted{5}4} (g) et \uline{Der \deleted{6}5} (g)

\deleted{\uline{Cor 2} (g) Soit \unsure{$S$|$\mathscr S$}
\diagram{\begin{tikzcd}[ampersand replacement=\&,row sep=small]
 X'\arrow[rr,"\alpha"]\arrow[dr,"f'"'] \& \& X\arrow[dl,"f"]\\
 \& S \&
\end{tikzcd}}
comm. dans Diag, avec $\alpha$ hot-équiv. stricte. Alors $\forall F$ dans $\mathbb D(X)$, on a
\[
 \alpha^*:f_*(F)\xrightarrow{\sim} f'_*(\alpha^*(F))
\]}

% slanted
\margin{Critère de changt de base pour images directes cohomologiques}
\uline{Cor. \unsure{3|5}} (g) Soit
\diagram{\begin{tikzcd}[ampersand replacement=\&,column sep=large]
 X'\arrow[r,"q"]\arrow[d,"f'"'] \& X\arrow[d,"f",""{name=right,coordinate}]\\
 Y'\arrow[r,"p",""{name=bottom,coordinate}] \& Y
 \arrow[from=right,to=bottom,bend right=12,"\unsure{\gamma|r}"']
\end{tikzcd}}
dans Diag (cf page 30), d'où hom de chgt de base
\[
 p^*f_*(F)\longrightarrow f'_*q^*F\qquad(F\in\mathrm{Ob}\,\mathbb D(X)).
\]
Pour que ce soit un iso. pour tout $F$, il suffit que \deleted{\unsure{le}} $\forall y'\in\mathrm{Ob}\,Y'$, le foncteur induit par $(p,q)$
\[
 X'_{/y'}\longrightarrow X_{/p(y')}
\]
une \uline{hot-équivalence stricte}\inserted{ (ou costricte)} ; cela signifie aussi
% slanted
\margin{\uline{Faux} pour ``stricte'', ne marche que pour ``\uline{co}stricte''}
\inserted{(\deleted{\unsure{dans}} \inserted{dans} le cas strict)} que pour
\[
 y'_0\in\mathrm{Ob}\,Y',\quad x_0\in\mathrm{Ob}\,X,\quad u_0:f(x_0)\to p(y'_0)
\]
la catégorie \inserted{$K(y'_0,x_0,u_0)$} des objets de la forme
\[
 (x',v',u)\mid x'\in\mathrm{Ob}\,X',\quad v':f'(x')\to y'_0,\quad u:q(x')\to x_0
\]
rendant commutatif le carré dans $Y'$
\diagram{$(*)$\quad\begin{tikzcd}[ampersand replacement=\&,column sep=large]
 f(q(x'))\arrow[r,"\gamma(x')"]\arrow[d,"f(u)"'] \& p(f'(x'))\arrow[d,"p(v')"]\\
 f(x_0)\arrow[r,"u_0"']\arrow[u,dashed,shift right=1ex,"\text{\scriptsize en cas de hoto-équiv. costricte}"'] \& p(y'_0)
\end{tikzcd}}
% slanted
\margin{comparer p. 30 \unsure{ff.}}
% slanted
\margin{\uline{NB} Dans le cas de la hot. équivalence \uline{costricte}, c'est la \unsure{=|même} description de $K={}_{(x_0,u_0)\backslash}(X'_{/\unsure{y'_0}})$, sauf que la flèche $u$ va en sens inv., $x_0\to q(x')$, et idem dans le diagramme \unsure{carré} $(*)$}
est \unsure{contractile|contractible}.

On applique \deleted{\unsure{Dér 5 (g)}} \inserted{Dér. 2, Der 4(g)} et \uline{Cor 1} \unsure{$\geq$|$\leq$|et}
% archive sheet number: 42
\scanpage{44}\deleted{\unsure{48|49}}\mspage{50}\deleted{\unsure{a|2}}
\deleted{Je vais généraliser légèrement le \deleted{\unsure{lemme}} \inserted{Cor. 1}}

\uline{Cor. 5} (g) \deleted{Soit} Soit $f:X\to Y$ pleinement fidèle dans Diag, et $F\in\mathbb D(X)$. Alors l'hom d'adjonction
\[
 f^*f_*(F)\longrightarrow F
\]
est un \uline{iso}.

En effet, il suffit de le vérifier fibre par fibre (Der 2). Soit, $x\in X$, \inserted{$y=f(x)\in Y$,} alors
\[
 f^*f_*(F)_x\simeq f_*(F)_y\underset{\text{Dér 5}}{\simeq}\deleted{\unsure{X_{/x}}}\Gamma_{\mathbb D}(X_{/y},F)
\]
or \deleted{\unsure{$X_{/y}$}} à cause de la pleine fidélité $X_{/y}\simeq X_{/x}$. \deleted{\unsure{Mais} Mais \unsure{j'ai|il} \unsure{oublié} l'\unsure{axiome}}

% slanted
\margin{Va résulter de \uline{Der 6} (g), plus fort ci-dessous}
% slanted
\margin{Axiome : \deleted{\unsure{limites}}\deleted{\unsure{coefficients}} la cat. d'indices à objet final (coeff. \uline{quelc.})}
\deleted{\unsure{Dér} \unsure{4}$_0$ (g)} : \inserted{\uline{Cor. 3} (\unsure{remonte|renvoi à} p. 48)} Soit $Z\in\mathrm{Diag}$, supposons que $Z$ a un objet final, \inserted{$e_Z$,} \deleted{plus généralement que le foncteur $\underleftarrow{\lim}_{Z}$ \unsure{ensembliste} (pour les préfaisceaux d'ens sur $Z$) transforme épi en épi (exactitude à droite de $\underleftarrow{\lim}$). Alors \unsure{Alors}} Soit $p:e\to Z$ le foncteur correspondant. Alors \deleted{\unsure{le}} l'hom canonique (pour $F$ dans $\mathbb D(Z)$)
\diagram{\begin{tikzcd}[ampersand replacement=\&,row sep=large]
 Z\arrow[d,"f"]\\ e\arrow[u,bend left=24,dashed,"p"]
\end{tikzcd}}
\[
 f_*(F)\to p^*(F)
\]
est un iso.

Cet hom \deleted{se déduit} est \unsure{le particulier|un particulier} de celui-ci, pour
% slanted
\margin{déjà envisagé dans Cor. 2 p. 47}
\diagram{\begin{tikzcd}[ampersand replacement=\&,row sep=small]
 Z'\arrow[rr,"p"]\arrow[dr,"f'"'] \& \& Z\arrow[dl,"f"]\\
 \& S \&
\end{tikzcd}}
(en faisant $Z'=S=e$, $f'=\mathrm{id}_e$)
\[
 f_*(F)\to f'_*(p^*(F))\qquad(F\in\mathrm{Ob}\,\mathbb D(Z))
\]
et c'est un cas particulier du morphisme de
% archive sheet number: 43
\scanpage{45}\mspage{51}
\unsure{changt|changement} de base pour
\diagram{\begin{tikzcd}[ampersand replacement=\&]
 Z'\arrow[r,"p"]\arrow[d,"f'"'] \& Z\rlap{$\;F$}\arrow[d,"f"]\\
 S\arrow[r,"\mathrm{id}"]\&S
\end{tikzcd}\quad $.\quad\big]$}

\uline{Remarque} : On retrouve le fait que $i^{\unsure{*}}i_*(F_Y)\xrightarrow{\sim}F_Y$ est iso, pour $i$ une imm. fermée. \deleted{\unsure{Retrouvons}} De plus avec les notations de Dér 4 (g), \deleted{\unsure{si on}} \deleted{\unsure{met}} $j^*(i_*(F_Y))$ se calcule par le th. de chgt de base (cor. 2 ci-dessus), qui s'applique dans la situation commutative
\diagram{\begin{tikzcd}[ampersand replacement=\&]
 X'\arrow[r,"q"]\arrow[d,"f'"']\arrow[dr,phantom,"\mathrm{cart}" description] \& X\arrow[d,"f"]\\
 Y'\arrow[r,"p"']\&Y
\end{tikzcd}}
quand $p$ est une imm. ouverte (et $\unsure{\gamma|r}=\mathrm{id}$) (alors $X'_{/y'}\to X_{/p(y')}$ est \unsure{$=$} un iso !), et on trouve, \unsure{revenant|venant} aux notations de Der 4 (g)
\diagram{\begin{tikzcd}[ampersand replacement=\&]
 \varnothing\arrow[r,"\unsure{q|g}"]\arrow[d,"f'"']\& Y\arrow[d,hook,"i"]\\
 U\arrow[r,hook,"j"']\& X
\end{tikzcd}}
ici : $j^*(i_{\deleted{\unsure{!|*}}*}\unsure{(F_Y)|[F_Y]}\simeq f'_*(\varphi)$, où $\varphi=\unsure{g|q}^*(F_Y)$ est l'``unique'' élément de la catégorie (équivalente à une cat. ponctuelle, par Der 1) $\mathbb D(\varnothing)$. D'autre part $f'_*$ transforme objet final en objet final, d'où $j^*i_*(F)$ est objet final.

Inversement, soit $F$ dans $\mathbb D(X)$, telle que $j^*(F)=e_{\mathbb D(U)}$. Considérons
\[
 F\xrightarrow{u}i_*i^*F\qquad(\mathrm{adj})
\]
je dis que c'est un \uline{iso}. En effet, par l'axiome
% archive sheet number: 44
\scanpage{46}\mspage{52}
de localisation Der 2, il suffit de prouver que $i^*(u)$ et $j^*(u)$ sont iso. Pour $i^*(u)$ cela résulte du fait que $i^*$ soit foncteur de localisation, et pour $j^*$ du \deleted{\unsure{l'h}} fait que $j^*(F)$ et $j^*(i_*i^*(F))$ sont des objets finaux dans $\mathbb D(U)$.

% slanted
\margin{\deleted{Axiome de comparaison par cofinalité cohomologique}}
\deleted{\uline{Dér 6} (g) : Soit $f:X\to Y$ \deleted{un morphisme} \inserted{dans Diag, tel} que $f$ soit une \unsure{hot-|hoto-}équivalence co-stricte (i.e. $f^\circ$ \unsure{hot|hoto} équiv. stricte). Alors $\forall F\in\mathbb D(Y)$, l'hom can
\[
 f^*:\Gamma_{\mathbb D}(Y,F)\to\Gamma_{\mathbb D}(X,f^*F)\tag*{(*)}
\]
est un iso.

\uline{NB} On a un énoncé plus fort si $f$ \inserted{elle-\unsure{=|même}} est une hot équivalence stricte (pas \unsure{néces.|néc.} costricte), cf Cor. 2 p. 47.

\deleted{\uline{\unsure{Question}}} On peut se demander de caractériser les
\[
 f:X\to Y
\]
de Diag, tels que $\forall$ dérivateur exact à gauche $\mathbb D$ (ou ne serait-ce que pour \deleted{\unsure{$\mathrm{HOTAb}$}} $\mathrm{Der}(\mathrm{Ab})$ et pour $\mathrm{HOT}$), et tt $F\in\mathbb D(Y)$, l'hom $(*)$ ci-dessus soit iso. \uline{Il n'est pas vrai} que $f$ soit une hoto-équivalence costricte. (soit \deleted{\unsure{coho}} cohomologiquement cofinal), car sinon une hot-équiv. costricte serait stricte, et l'inv. serait vrai par dualité, mais ce n'est pas vrai.}
% archive sheet number: 45
\scanpage{47}\mspage{53}
\uline{Morphismes \deleted{\unsure{co}}homotopiquement propres, homotopiquement lisses}

Soit $f:X\to Y$ morphisme de catégories, considérons $y_0\in\mathrm{Ob}\,\unsure{Y_0|Y}$, et le foncteur canonique
\[
 \underset{\substack{\text{fibre de }X\\\text{en }y_0}}{X_{y_0}}\xrightarrow{\varphi}X_{/y_0}
 \qquad \varphi(x)=(x,\mathrm{id}_{f(x)})\tag*{(1)}
\]
Soit $(x_0,u_0)\in\mathrm{Ob}\,X_{/y_0}$ donc
\[
 \begin{aligned}x_0&\in\mathrm{Ob}\,X\\u_0&:f(x_0)\to y_0\end{aligned}
\]
considérons la catégorie coinduite par $\varphi$ sur $(x_0,u_0)$
\[
 {}_{(x_0,u_0)\backslash}(X_{y_0})\deleted{\unsure{/(x_0,u_0)}},
\]
pour tester si $\varphi$ est une \deleted{\unsure{éq}} \unsure{hotéquivalence|hotoéquivalence} \inserted{co}stricte\inserted{ i.e. \unsure{homol.|homol.g}$^{\mathrm t}$ cofinal}. \deleted{On} On va l'interpréter via le \deleted{morphisme} \inserted{foncteur} induit par $f$
\[
 {}_{x_0\backslash}X\xrightarrow{\psi}{}_{f(x_0)\backslash}Y
\]
allant de la cat. des spécialisations de $x_0$ \inserted{de $X$,} vers celle des spécialisations de $f(x_0)$ dans $Y$. On trouve alors un iso can. \unsure{avec|entre} la fibre de $\psi$ en $u_0$

\diagram{\begin{tikzcd}[ampersand replacement=\&,row sep=small]
 x_0\arrow[r,dashed]\&x\\
 f(x_0)\arrow[r,"u_0"]\&y_0
\end{tikzcd}}
\[
 \begin{aligned}
 &\text{Catégorie coinduite }{}_{(x_0,u_0)\backslash} X_{y_0}\text{ par }X_{y_0}\text{ au dessus}\\
 &\text{de \unsure{de} l'objet }(x_0,u_0)\text{ de }X_{/y_0}\\
 &\simeq\text{ Catégorie fibre }({}_{x_0\backslash}X)_{u_0}\text{ de }{}_{x_0\backslash}X\text{ au dessus}\\
 &\text{de } {}_{f(x_0)\backslash}Y\text{ en }u_0
 \end{aligned}
\]
Donc

\uline{Proposition} Pour que les foncteurs $\varphi:X_{y_0}\to X_{/y_0}$ \inserted{($y_0\in\mathrm{Ob}\,Y$)} soient des hotéquiv. costrictes, il f. et s. que les fibres des foncteurs ${}_{x_0\backslash}X\longrightarrow {}_{f(x_0)\backslash}Y$ induits par $f$ soient contractiles.
% archive sheet number: 46
\scanpage{48}\mspage{54}
On dit alors que $f$ est \uline{\deleted{\unsure{co}}homotopiquement propre}. La \deleted{\unsure{condition}} \inserted{notion} duale s'exprime, soit par la condition que les foncteurs induits par $f$
\[
 X_{y_0}\longrightarrow {}_{y_0\backslash}Y
\]
soient des hotéquivalences strictes i.e. qu'ils soient cohomotopiquement cofinaux, soit que les fibres des morphismes induits sur les ``catégories localisées'' (catégories des générisations)
\[
 X_{/x_0}\longrightarrow Y_{/f(x_0)}
\]
sont contractiles. (N.B., On utilise ici qu'une catégorie $Z$ est contractile ssi $Z^\circ$ l'est.) On dit alors que $f$ est \uline{homotopiquement lisse}

\uline{Prop} a) La notion de foncteur propre \inserted{(\unsure{resp.|ou} lisse)} est stable par changement de base, par composition. \inserted{Les isomorphismes.} \deleted{Plus gén. les \unsure{si} $gf$ est propre, et $g$ propre, alors $f$ est propre ; i.e. tout $S$-foncteur entre catégories $X,Y$ sur $S$ propres sur $S$, est propre.}

% slanted
\margin{Attention, une équiv. de cat. n'est pas propre !}
\deleted{b) L'énoncé dual pour foncteurs lisses. Notamment, si $gf$ est lisse et \unsure{l'équivalence de catégories}}
\inserted{\unsure{Mais oui|Mais} les immers. fermées, sont propres. L'image d'un morphisme propre est une \deleted{\unsure{partie de}} sous-catégorie fermée (\unsure{non pleine})$^*$}
% slanted
\margin{Si $X'\hookrightarrow X$ inclusion d'une sous-catégorie de $X$, elle \deleted{\unsure{est}} est propre \unsure{ssi|$\Leftrightarrow$} \deleted{\unsure{si}} c'est une \unsure{immers.|immersion} fermée}

\deleted{\unsure{Cor.}} \uline{NB} le foncteur $X\to e$ est tjrs propre et lisse, car les cat ${}_{x_0\backslash}X$ et $X_{/x_0}$ sont contractiles (ayant objet initial ou objet final). Donc tt foncteur de projection $X\times Y\xrightarrow[\mathrm{pr}_{\unsure{1|,}}]{}Y$ est à la fois propre et lisse.
% archive sheet number: 47
\scanpage{49}\mspage{55}
\uline{Dém} Stabilité par changement de base :
\diagram{\begin{tikzcd}[ampersand replacement=\&]
 X \arrow[d,"f"'] \& X' \arrow[l,"q"'] \arrow[d,"f'"] \arrow[dl,phantom,"\mathrm{cart}" description]\\
 Y \& Y' \arrow[l,"p"]
\end{tikzcd}}
d'où, pour $x'\in\mathrm{Ob}X'$
\diagram{\begin{tikzcd}[ampersand replacement=\&]
 X_{/x} \arrow[d] \& X'_{/x'} \arrow[l] \arrow[d]\\
 Y_{/y} \& Y'_{/y'} \arrow[l]
\end{tikzcd}}
où $x,y',y$ sont les images de \unsure{$x$|$x'$} dans $X,Y',Y$.

On vérifie que ce diagramme est \uline{cartésien}, donc les fibres de $X'_{/x'}\longrightarrow Y'_{/y'}$ sont \unsure{isom.} à des fibres de $X_{/x}\to Y_{/y}$ donc contractiles si celles-ci le sont.

Dans \unsure{un|notre} cas $X\xrightarrow{f}Y$ on vérifie que \deleted{les bilocalisés pour} \inserted{pour un foncteur} \deleted{\unsure{lisse un foncteur} \unsure{ou de}} bilocalisé de $f$
\[
 X_{/x}\xrightarrow{\varphi\unsure{=|\times}f_{/x}}Y_{/y},
\]
\deleted{\unsure{s'il}} \deleted{\unsure{tout}} ``bilocalisé''
\[
 (X_{/x})_{/\xi}\longrightarrow(Y_{/y})_{/\eta}
\]
\inserted{$(\xi:x'\to x,\eta:y'\to y)$} de $\varphi$ s'identifie\unsure{nt} \unsure{encore} à \unsure{un} bilocalisé de $f$
\[
 X_{/x'}\longrightarrow Y_{/y'}
\]
Donc si $f$ est \deleted{propre} \inserted{lisse}, les bilocalisés $f_{/x}$ $(x\in\mathrm{Ob}X)$ le sont aussi. Notons d'autre part

\uline{Proposition} Tout \deleted{\unsure{foncteur}} propre \inserted{(on dirait lisse)} à fibres contractiles est une hot-équivalence.

% slanted
\margin{Cf. \unsure{utilisation|indication|imbrication} \unsure{ds cor. 2|des car. 2} \unsure{pour|par} \illegible{1}}
% archive sheet number: 48
